\chapter{Discusiones} \label{discusiones}

\section{Estudio sobre la influencia de marcadores en el declive cognitivo en la Enfermedad del Parkinson}
En los resultados obtenidos a través del algoritmo de análisis en Python se proporciona información crucial sobre la influencia de las variables estudiadas en el deterioro cognitivo en pacientes con Enfermedad de Parkinson.\\

Observando las dos tablas resultantes, se puede llevar a algunas conclusiones sobre los distintos parámetros. En la tabla \ref{tab:tabla1}, tras la comparación de los valores en la linea de base, se observan diferencias relevantes en algunas filas entre el grupo de pacientes con Parkinson y el grupo de control. En primer lugar, cabe destacar que para que el estudio sea fiable, ambos grupos tiene las mismas características demográficas. La edad, a pesar de lo que se podría llegar a pensar ya que la media de edad de ambos grupos es bastante elevada, no es un indicador del avance hacia el deterioro. Sin embargo, entre los parámetros cognitivos y clínicos se encuentran varios marcadores con indicador de significancia en la columna del Post-hoc. Entre los que se pueden señalar están GDS, REM, MOCA, SDMT y MEANPUT. \\

Por otra parte, en la tabla \ref{tab:tabla2} destaca la relación entre las distintas características de estudio y el desarrollo de diferentes niveles de declive cognitivo en los pacientes PD. Los valores de p-valor y los indicadores de post-hoc señalan que variables pueden tener un efecto significativo en el transcurso del deterioro cognitivo en los distintos subgrupos. Entre ellas, diferenciadas por tener un \textit{valor p} igual o por debajo de 0.001, se ve como por ejemplo el SDMT y el MEANCAU son considerables tanto en la comparación entre sujetos cognitivamente normales y con deterioro cognitivo como entre cognitivamente normales y los pacientes que desarrollaron demencia. Sin embargo, hay factores que solo muestran efectos significativos en una de las dos comparativas, marcados con un solo indicador (\textit{a} o \textit{b}). Entre estos últimos se encuentran REM, con diferencias importantes solo entre PD\_CU y PD\_MCI; o A$\beta_{42}$, con diferencia notable entre PD\_CU y PDD. En esta tabla, dentro de las variables demográficas, la edad es un indicador del avance de los pacientes hacia un deterioro cognitivo medio y de la aparición de la demencia. Además, se puede ver como las variables cognitivas tienen una gran importancia en el desarrollo del declive cognitivo, sobre todo indicando la evolución de PD\_CU a PD\_MCI. Aún así, los parámetro LNS, SDMT, INMEDIATE y SFT también marcan la evolución de los sujetos hacia el grupo con demencia. Entre las variables clínicas, ABETA y PATU señalan la evolución hacia el colectivo PDD, y MEANCAU es una característica significativa tanto de la aparición del deterioro cognitivo como de la demencia. \\

En vistas generales, cabe destacar las variables que resultaron significativas en el análisis, tanto en la tabla \ref{tab:tabla1} como en la tabla \ref{tab:tabla2}. Entre estas, se encuentran: UPSIT, SDMT, INMEDIATE, SFT y MEANCAU. Todas ellas son medidas cognitivas y clínicas. Estos parámetros pueden tener implicaciones cruciales para el diagnostico anticipado y el desarrollo de estrategias para el tratamiento.\\

En resumen, los resultados de este estudio destacan marcadores significativos en el declive cognitivo en la Enfermedad de Parkinson. Estos hallazgos tienen el potencial de mejorar la predicción temprana de la enfermedad y permitir intervenciones más efectivas. La identificación de estos marcadores puede conducir a un enfoque más personalizado en el tratamiento de pacientes con Parkinson, mejorando así su calidad de vida y ofreciendo esperanza a quienes se ven afectados por esta enfermedad. También es importante tener en cuenta que los resultados obtenidos pueden variar según la población y otros factores.


\section{Análisis longitudinal y reconstrucción temporal de un único sujeto}

En las gráficas que podemos ver en las figuras \ref{fig:conv1_gvsource}, \ref{fig:conv1_tauxi} y \ref{fig:conv1_noise} se demuestra la convergencia de los parámetros estimados por Leaspy a lo largo de 50 iteraciones. Indicando la validez de los parámetros (g, v0, sources, tau, xi, noise).\\

En el caso de las gráficas obtenidas que estiman las trayectorias \textit{'short-term'} de cada variable, se puede observar en la figura \ref{fig:trayM5norm} como en los parámetros ABETA, MEANCAU Y HVLTRDLY su estimación no se cruza con la curva de unión de los puntos de sus observaciones. En las otras tres (SDMT, LNS, SFT) si que se puede ver una clara intersección entre las rectas. Esta diferencia se puede deber al número de visitas del paciente, ya que a mayor numero de mediciones mas precisa es la estimación.


\section{Modelos de progresión con múltiples sujetos}
En primer lugar, se obtuvieron las tablas de medición del NMAE para los diez bootstraps del grupo de individuos de tipo \textit{test} (tabla \ref{tab:NMAE test}) y de tipo \textit{train} (tabla \ref{tab:NMAE train}). Luego se obtuvieron los mismo valores en ecala originalr, representados en las tablas \ref{tab:mae test} y \ref{tab:mae train}. en el caso de los valores NMAE se obtuvieron valores cercanos a cero, un indicador positivo de la estimación realizada. En cuanto a los MAEs los resultados obtenidos reflejan valores cercanos a los valores originales de los biomarcadores también indicando buena precisión.  \\

Como es lógico en las prestaciones para el grupo de entrenamiento se obtienen resultados mas exactos que con el grupo de test, ya que los parámetros calculados en el método de personalización se llevo a cabo con el primer grupo nombrado.\\

La siguiente tabla que se analiza corresponde a la figura \ref{fig:SCORES}. El SCORE1 tanto en \textit{test} como en \textit{train} son aparentemente correcto y concordante con los visto anteriormente. Esto indicaría que se ha hecho una estimación con una precisión ligeramente por encima del 50\% del número de pacientes cogniticamente normales al inicio de las observaciones.
El resto de SCORES muestran valores demasiado altos. Se puede asumir que el cálculo de estas prestaciones es incorrecto. 